\documentclass[12pt, a4paper]{article}

\usepackage{xcolor}
\usepackage{geometry}
\usepackage{hyperref}
\usepackage{titlesec}
\usepackage{enumitem}
\usepackage{longtable}
\usepackage{booktabs}
\usepackage{array}
\usepackage{ragged2e}
\usepackage{fancyhdr}
\usepackage{tcolorbox}
\usepackage{amssymb}
\usepackage{listings}

\geometry{margin=2.5cm}

\usepackage{xepersian}
\settextfont{Vazirmatn}
\setlatintextfont{Times New Roman}
\setdigitfont{Vazirmatn}

\hypersetup{
    colorlinks=true,
    linkcolor=blue,
    urlcolor=blue
}

\titleformat{\section}{\Large\bfseries\color{blue!60!black}}{\thesection}{1em}{}
\titleformat{\subsection}{\large\bfseries}{\thesubsection}{1em}{}
\titleformat{\subsubsection}{\normalsize\bfseries}{\thesubsubsection}{1em}{}

\setlist{nosep, leftmargin=*}

\pagestyle{fancy}
\fancyhf{}
\fancyhead[R]{\small جزوهٔ مطالعه — جلسهٔ ۷}
\fancyhead[L]{\small الگوهای طراحی}
\fancyfoot[C]{\thepage}
\renewcommand{\headrulewidth}{0.4pt}

\linespread{1.4}

\lstset{
    basicstyle=\ttfamily\small,
    breaklines=true,
    frame=single,
    language=Java,
    showstringspaces=false
}

\title{
    \vspace{-1cm}
    \textbf{\Huge جزوهٔ مطالعه}\\[0.3cm]
    \textbf{\LARGE جلسهٔ ۷: الگوهای طراحی}\\[0.5cm]
    \large مهندسی نرم‌افزار
}
\author{
    تهیه‌کننده: ابراهیم فضلی\\
    ترم: ۱۴۰۴
}
\date{\today}

\begin{document}

\maketitle
\thispagestyle{empty}

\vspace{1cm}

\begin{center}
\begin{tabular}{|l|l|}
\hline
\textbf{درس} & مهندسی نرم‌افزار \\
\hline
\textbf{جلسه} & ۷ \\
\hline
\textbf{موضوع} & Design Patterns (Factory, Strategy, Observer, ...) \\
\hline
\textbf{زمان مطالعه} & ۳-۴ ساعت \\
\hline
\textbf{پیش‌نیاز} & جلسات ۱-۶ \\
\hline
\end{tabular}
\end{center}

\newpage
\tableofcontents
\newpage

% ============================================================
\section{مقدمه}

\subsection{چرا Design Pattern؟}

در جلسهٔ ۶ یاد گرفتیم:
\begin{itemize}
    \item چهار اصل OOP
    \item SOLID
    \item Coupling و Cohesion
\end{itemize}

حالا سؤال: \textbf{چطور این اصول را در عمل پیاده کنیم؟}

\textbf{پاسخ:} Design Patterns.

\subsection{Design Pattern چیست؟}

\begin{quote}
\textbf{Design Pattern} یک راه‌حل آماده برای یک مشکل رایج طراحی است.
\end{quote}

\textbf{نکته مهم:} الگو = \textbf{قالب}، نه کد آماده.

\subsection{مثال ساده}

\textbf{مشکل:} ساخت روش‌های پرداخت مختلف (کارت، PayPal، کریپتو).

\textbf{راه‌حل بدون الگو:}
\begin{lstlisting}
public void pay(String method, double amount) {
    if (method.equals("card")) { ... }
    else if (method.equals("paypal")) { ... }
    else if (method.equals("crypto")) { ... }
    // هر روش جدید = تغییر این متد
}
\end{lstlisting}

\textbf{راه‌حل با Factory:}
\begin{lstlisting}
PaymentMethod method = PaymentFactory.create("card");
method.pay(100000);
// افزودن روش جدید = فقط کلاس جدید
\end{lstlisting}

\textbf{مزیت:} کد باز برای توسعه، بسته برای تغییر (OCP).

\subsection{آنچه در این جزوه یاد می‌گیرید}
\begin{itemize}
    \item سه دسته Design Pattern
    \item Creational: Factory, Singleton, Builder
    \item Structural: Adapter, Decorator, Facade
    \item Behavioral: Strategy, Observer
    \item Repository Pattern
    \item مثال‌های HomeStore
\end{itemize}

% ============================================================
\section{سه دسته Design Pattern}

\begin{center}
\begin{tabular}{|l|p{9cm}|}
\hline
\textbf{دسته} & \textbf{هدف} \\
\hline
Creational & ساخت اشیاء \\
\hline
Structural & ترکیب اشیاء \\
\hline
Behavioral & تعامل اشیاء \\
\hline
\end{tabular}
\end{center}

\subsection{Creational Patterns}
\begin{itemize}
    \item \textbf{Factory Method} — ساخت اشیاء با interface مشترک
    \item \textbf{Abstract Factory} — ساخت خانواده‌ای از اشیاء
    \item \textbf{Singleton} — فقط یک نمونه
    \item \textbf{Builder} — ساخت گام‌به‌گام
    \item \textbf{Prototype} — کپی از نمونه
\end{itemize}

\subsection{Structural Patterns}
\begin{itemize}
    \item \textbf{Adapter} — تبدیل interface
    \item \textbf{Decorator} — افزودن رفتار
    \item \textbf{Facade} — ساده‌سازی
    \item \textbf{Proxy} — کنترل دسترسی
    \item \textbf{Composite} — ساختار درختی
\end{itemize}

\subsection{Behavioral Patterns}
\begin{itemize}
    \item \textbf{Strategy} — الگوریتم قابل تعویض
    \item \textbf{Observer} — اطلاع‌رسانی
    \item \textbf{Command} — کپسوله‌سازی درخواست
    \item \textbf{State} — تغییر رفتار با حالت
    \item \textbf{Template Method} — قالب الگوریتم
\end{itemize}

% ============================================================
\section{Creational Patterns}

\subsection{Factory Method}

\subsubsection{مسئله}
کلاینت نمی‌داند کدام کلاس را بسازد. تصمیم ساخت باید در یک جا باشد.

\subsubsection{مثال HomeStore}
\begin{lstlisting}
public interface PaymentMethod {
    void pay(double amount);
}

public class CardPayment implements PaymentMethod {
    public void pay(double amount) { ... }
}

public class PayPalPayment implements PaymentMethod {
    public void pay(double amount) { ... }
}

public class PaymentFactory {
    public static PaymentMethod create(String type) {
        if (type.equals("card")) return new CardPayment();
        if (type.equals("paypal")) return new PayPalPayment();
        throw new IllegalArgumentException();
    }
}

// استفاده
PaymentMethod payment = PaymentFactory.create("card");
payment.pay(100000);
\end{lstlisting}

\subsubsection{مزیت}
\begin{itemize}
    \item افزودن روش جدید = فقط کلاس جدید
    \item کلاینت از جزئیات مستقل است
    \item تست آسان‌تر
\end{itemize}

\subsection{Abstract Factory}

\subsubsection{تعریف}
ساخت \textbf{خانواده‌ای} از اشیاء مرتبط.

\subsubsection{مثال}
\begin{lstlisting}
public interface IUIFactory {
    Button createButton();
    Form createForm();
    Table createTable();
}

public class LightThemeFactory implements IUIFactory {
    public Button createButton() { return new LightButton(); }
    // ...
}

public class DarkThemeFactory implements IUIFactory {
    public Button createButton() { return new DarkButton(); }
    // ...
}
\end{lstlisting}

\subsubsection{تفاوت با Factory Method}
\begin{itemize}
    \item \textbf{Factory Method:} یک شیء
    \item \textbf{Abstract Factory:} خانواده‌ای از اشیاء
\end{itemize}

\subsection{Singleton}

\subsubsection{تعریف}
فقط یک نمونه از کلاس.

\subsubsection{مثال}
\begin{lstlisting}
public class DatabaseConnection {
    private static DatabaseConnection instance;
    
    private DatabaseConnection() { ... }
    
    public static DatabaseConnection getInstance() {
        if (instance == null) {
            instance = new DatabaseConnection();
        }
        return instance;
    }
}
\end{lstlisting}

\subsubsection{کاربرد}
\begin{itemize}
    \item اتصال به دیتابیس
    \item تنظیمات برنامه
    \item لاگ کردن
\end{itemize}

\subsubsection{هشدار}
\begin{itemize}
    \item تست سخت می‌شود
    \item Coupling بالا
    \item استفاده با احتیاط
\end{itemize}

\subsection{Builder}

\subsubsection{تعریف}
ساخت گام‌به‌گام اشیاء پیچیده.

\subsubsection{مثال}
\begin{lstlisting}
public class OrderBuilder {
    private Order order = new Order();
    
    public OrderBuilder addItem(OrderItem item) {
        order.addItem(item);
        return this;
    }
    
    public OrderBuilder setAddress(String address) {
        order.setAddress(address);
        return this;
    }
    
    public OrderBuilder setPayment(PaymentMethod payment) {
        order.setPayment(payment);
        return this;
    }
    
    public Order build() {
        return order;
    }
}

// استفاده
Order order = new OrderBuilder()
    .addItem(item1)
    .addItem(item2)
    .setAddress("Tehran")
    .setPayment(new CardPayment())
    .build();
\end{lstlisting}

\subsubsection{مزیت}
کد خواناتر از constructor با ۱۰ پارامتر.

% ============================================================
\section{Structural Patterns}

\subsection{Adapter}

\subsubsection{تعریف}
تبدیل interface ناسازگار به interface مورد انتظار.

\subsubsection{مثال}
\begin{lstlisting}
// کد ما
public interface PaymentGateway {
    void pay(double amount);
}

// درگاه خارجی
public class StripeAPI {
    public void charge(double amount) { ... }
}

// Adapter
public class StripeAdapter implements PaymentGateway {
    private StripeAPI stripe;
    
    public StripeAdapter() {
        this.stripe = new StripeAPI();
    }
    
    public void pay(double amount) {
        stripe.charge(amount);
    }
}
\end{lstlisting}

\subsubsection{مزیت}
کد ما از جزئیات درگاه مستقل می‌شود.

\subsection{Decorator}

\subsubsection{تعریف}
افزودن رفتار بدون تغییر کلاس اصلی.

\subsubsection{مثال}
\begin{lstlisting}
public interface OrderService {
    void process(Order order);
}

public class OrderServiceImpl implements OrderService {
    public void process(Order order) { ... }
}

public class LoggedOrderService implements OrderService {
    private OrderService inner;
    
    public LoggedOrderService(OrderService inner) {
        this.inner = inner;
    }
    
    public void process(Order order) {
        System.out.println("Processing: " + order);
        inner.process(order);
    }
}

// استفاده
OrderService service = new LoggedOrderService(
    new OrderServiceImpl()
);
\end{lstlisting}

\subsubsection{مزیت}
ترکیب رفتارها بدون تغییر کد اصلی.

\subsection{Facade}

\subsubsection{تعریف}
ساده‌سازی interface پیچیده.

\subsubsection{مثال}
\begin{lstlisting}
public class PurchaseFacade {
    private OrderService orderService;
    private PaymentService paymentService;
    private EmailService emailService;
    private StockService stockService;
    
    public void buy(Order order, PaymentMethod payment) {
        orderService.create(order);
        paymentService.charge(payment, order.getTotal());
        stockService.reserve(order.getItems());
        emailService.sendConfirmation(order);
    }
}

// استفاده - فقط یک متد
PurchaseFacade facade = new PurchaseFacade();
facade.buy(order, payment);
\end{lstlisting}

\subsubsection{مزیت}
کلاینت فقط یک متد صدا می‌زند.

% ============================================================
\section{Behavioral Patterns}

\subsection{Strategy}

\subsubsection{مسئله}
محاسبهٔ هزینهٔ ارسال:
\begin{itemize}
    \item پست: بر اساس وزن
    \item تیپاکس: بر اساس شهر
    \item پیک: بر اساس فاصله
\end{itemize}

\subsubsection{راه‌حل}
\begin{lstlisting}
public interface ShippingStrategy {
    double calculate(double weight);
}

public class PostShipping implements ShippingStrategy {
    public double calculate(double weight) {
        return weight * 5000;
    }
}

public class ExpressShipping implements ShippingStrategy {
    public double calculate(double weight) {
        return weight * 15000;
    }
}

// استفاده
ShippingStrategy strategy = new ExpressShipping();
double cost = strategy.calculate(2.5);
\end{lstlisting}

\subsubsection{مزیت}
افزودن روش جدید = فقط کلاس جدید.

\subsection{Observer}

\subsubsection{مسئله}
وقتی سفارش ثبت می‌شود، چه کسی باید خبردار شود؟
\begin{itemize}
    \item ایمیل به مشتری
    \item SMS به انبار
    \item اطلاع به پنل ادمین
\end{itemize}

\subsubsection{راه‌حل}
\begin{lstlisting}
public interface OrderObserver {
    void onOrderCreated(Order order);
}

public class EmailNotifier implements OrderObserver {
    public void onOrderCreated(Order order) {
        // ارسال ایمیل
    }
}

public class SmsNotifier implements OrderObserver {
    public void onOrderCreated(Order order) {
        // ارسال SMS
    }
}

public class Order {
    private List<OrderObserver> observers = new ArrayList<>();
    
    public void addObserver(OrderObserver observer) {
        observers.add(observer);
    }
    
    public void create() {
        // ثبت سفارش
        for (OrderObserver observer : observers) {
            observer.onOrderCreated(this);
        }
    }
}
\end{lstlisting}

\subsubsection{مزیت}
افزودن notifier جدید = فقط کلاس جدید.

% ============================================================
\section{Repository Pattern}

\subsection{تعریف}
جداسازی منطق دیتابیس از منطق کسب‌وکار.

\subsection{مثال}
\begin{lstlisting}
public interface IOrderRepository {
    Order findById(int id);
    void save(Order order);
    List<Order> findByCustomer(int customerId);
}

public class PostgresOrderRepository implements IOrderRepository {
    public Order findById(int id) { ... }
    public void save(Order order) { ... }
    public List<Order> findByCustomer(int customerId) { ... }
}

public class InMemoryOrderRepository implements IOrderRepository {
    private Map<Integer, Order> orders = new HashMap<>();
    
    public Order findById(int id) { return orders.get(id); }
    public void save(Order order) { orders.put(order.getId(), order); }
    public List<Order> findByCustomer(int customerId) { ... }
}
\end{lstlisting}

\subsection{مزیت}
\begin{itemize}
    \item تغییر دیتابیس = فقط تغییر implementation
    \item تست آسان با InMemory
    \item منطق کسب‌وکار از دیتابیس مستقل
\end{itemize}

% ============================================================
\section{سؤالات کلیدی برای مطالعه}

\subsection{سطح ۱ (پایه)}

\subsubsection{سؤال ۱: Design Pattern چیست؟}
\textbf{پاسخ:} راه‌حل آماده برای مشکلات رایج طراحی. الگو = قالب، نه کد 
آماده.

\subsubsection{سؤال ۲: سه دسته Design Pattern چیست؟}
\textbf{پاسخ:}
\begin{itemize}
    \item \textbf{Creational:} ساخت اشیاء
    \item \textbf{Structural:} ترکیب اشیاء
    \item \textbf{Behavioral:} تعامل اشیاء
\end{itemize}

\subsubsection{سؤال ۳: Factory Method چیست؟}
\textbf{پاسخ:} یک interface مشترک + یک Factory که شیء مناسب را می‌سازد. 
کلاینت نمی‌داند کدام کلاس ساخته می‌شود.

\subsubsection{سؤال ۴: Strategy چیست؟}
\textbf{پاسخ:} الگوریتم قابل تعویض. یک interface + چند implementation. 
مثال: محاسبهٔ هزینهٔ ارسال با روش‌های مختلف.

\subsubsection{سؤال ۵: Observer چیست؟}
\textbf{پاسخ:} اطلاع‌رسانی تغییرات. یک Subject + چند Observer. مثال: وقتی 
سفارش ثبت می‌شود، ایمیل و SMS فرستاده می‌شود.

\subsection{سطح ۲ (متوسط)}

\subsubsection{سؤال ۶: تفاوت Factory Method و Abstract Factory چیست؟}
\textbf{پاسخ:}
\begin{itemize}
    \item \textbf{Factory Method:} یک شیء
    \item \textbf{Abstract Factory:} خانواده‌ای از اشیاء مرتبط
\end{itemize}

\subsubsection{سؤال ۷: چرا Singleton خطرناک است؟}
\textbf{پاسخ:}
\begin{itemize}
    \item تست سخت می‌شود (state مشترک)
    \item Coupling بالا
    \item چند رشته‌ای سخت
    \item جایگزین: Dependency Injection
\end{itemize}

\subsubsection{سؤال ۸: تفاوت Adapter و Decorator چیست؟}
\textbf{پاسخ:}
\begin{itemize}
    \item \textbf{Adapter:} تبدیل interface (تغییر ظاهر)
    \item \textbf{Decorator:} افزودن رفتار (تغییر عملکرد)
\end{itemize}

\subsubsection{سؤال ۹: Repository Pattern چه مزیتی دارد؟}
\textbf{پاسخ:}
\begin{itemize}
    \item جداسازی منطق دیتابیس از کسب‌وکار
    \item تست آسان با InMemory
    \item تغییر دیتابیس بدون تغییر کسب‌وکار
\end{itemize}

\subsubsection{سؤال ۱۰: چطور Design Pattern مناسب را انتخاب کنیم؟}
\textbf{پاسخ:}
\begin{itemize}
    \item مشکل را دقیق بشناس
    \item ببین کدام الگو با مشکل همخوانی دارد
    \item الگو را با نیاز انتخاب کن، نه از روی مد
    \item ساده‌ترین راه‌حل را انتخاب کن
\end{itemize}

\subsection{سطح ۳ (پیشرفته)}

\subsubsection{سؤال ۱۱: آیا همیشه باید Design Pattern استفاده کنیم؟}
\textbf{پاسخ:} نه. برای پروژه‌های کوچک، ممکن است پیچیدگی اضافه کند. Design 
Pattern برای پروژه‌هایی است که \textbf{طول عمر} دارند و \textbf{تغییر} 
می‌کنند.

\subsubsection{سؤال ۱۲: چطور Design Pattern را در تیم معرفی کنیم؟}
\textbf{پاسخ:}
\begin{itemize}
    \item آموزش
    \item مثال‌های خوب در کد
    \item Code Review با چک‌لیست
    \item مستندسازی تصمیم‌ها
\end{itemize}

\subsubsection{سؤال ۱۳: تفاوت Strategy و State چیست؟}
\textbf{پاسخ:}
\begin{itemize}
    \item \textbf{Strategy:} الگوریتم قابل تعویض، انتخاب توسط کلاینت
    \item \textbf{State:} تغییر رفتار با حالت، انتخاب خودکار
\end{itemize}

\subsubsection{سؤال ۱۴: چطور Design Pattern را تست کنیم؟}
\textbf{پاسخ:}
\begin{itemize}
    \item Unit Test برای هر implementation
    \item Mock برای interface
    \item Integration Test برای ترکیب
\end{itemize}

\subsubsection{سؤال ۱۵: اگر کد قدیمی Design Pattern را نقض می‌کند، چه کار 
کنیم؟}
\textbf{پاسخ:}
\begin{itemize}
    \item Refactor تدریجی (نه یک‌جا)
    \item اولویت‌بندی بر اساس تغییرات آینده
    \item تست قبل از Refactor
    \item هر Sprint یک الگو
\end{itemize}

% ============================================================
\section{تمرین‌ها}

\subsection{تمرین ۱: انتخاب الگو}
برای هر سناریو، الگوی مناسب را انتخاب کنید:
\begin{enumerate}
    \item محاسبهٔ مالیات برای کشورهای مختلف → Strategy
    \item ارسال ایمیل و SMS هنگام ثبت سفارش → Observer
    \item اتصال به درگاه‌های پرداخت مختلف → Adapter
    \item کش کردن نتایج جست‌وجو → Decorator
    \item ساخت گزارش با گزینه‌های زیاد → Builder
\end{enumerate}

\subsection{تمرین ۲: پیاده‌سازی Factory}
یک \lr{NotificationFactory} بنویسید که ۳ نوع اطلاع‌رسانی (ایمیل، SMS، Push) 
را پشتیبانی کند.

\subsection{تمرین ۳: پیاده‌سازی Strategy}
یک \lr{TaxStrategy} بنویسید که مالیات را برای ۳ کشور (ایران، ترکیه، آلمان) 
محاسبه کند.

\subsection{تمرین ۴: پیاده‌سازی Observer}
یک \lr{StockObserver} بنویسید که وقتی موجودی محصول کم می‌شود، به ادمین 
اطلاع دهد.

\subsection{تمرین ۵: Repository}
یک \lr{ProductRepository} بنویسید که هم PostgreSQL و هم InMemory را 
پشتیبانی کند.

% ============================================================
\section{منابع}

\subsection{منابع اصلی}
\begin{center}
\begin{tabular}{|l|l|}
\hline
\textbf{منبع} & \textbf{فصل} \\
\hline
Design Patterns — GoF & کل کتاب \\
\hline
Head First Design Patterns & فصل ۱-۵ \\
\hline
Clean Architecture & فصل ۱۱-۱۴ \\
\hline
\end{tabular}
\end{center}

\subsection{منابع آنلاین}
\begin{itemize}
    \item \textbf{ویدیو:} Design Patterns in 1 hour
    \item \textbf{وب‌سایت:} Refactoring.Guru
    \item \textbf{کتاب رایگان:} Design Patterns (GoF)
\end{itemize}

% ============================================================
\section{چک‌لیست مطالعه}

\subsection{قبل از مطالعه}
\begin{itemize}
    \item[$\square$] جزوه را دانلود کردم
    \item[$\square$] وقت کافی دارم (۳-۴ ساعت)
\end{itemize}

\subsection{در حین مطالعه}
\begin{itemize}
    \item[$\square$] بخش ۱ (مقدمه) را خواندم
    \item[$\square$] بخش ۲ (سه دسته) را خواندم
    \item[$\square$] بخش ۳ (Creational) را خواندم
    \item[$\square$] بخش ۴ (Structural) را خواندم
    \item[$\square$] بخش ۵ (Behavioral) را خواندم
    \item[$\square$] بخش ۶ (Repository) را خواندم
    \item[$\square$] بخش ۷ (سؤالات) را خواندم
    \item[$\square$] بخش ۸ (تمرین‌ها) را خواندم
\end{itemize}

\subsection{بعد از مطالعه}
\begin{itemize}
    \item[$\square$] سؤالات سطح ۱ را جواب دادم
    \item[$\square$] سؤالات سطح ۲ را جواب دادم
    \item[$\square$] سؤالات سطح ۳ را جواب دادم
    \item[$\square$] تمرین ۱ را انجام دادم
    \item[$\square$] تمرین ۲ را انجام دادم
    \item[$\square$] تمرین ۳ را انجام دادم
    \item[$\square$] تمرین ۴ را انجام دادم
    \item[$\square$] تمرین ۵ را انجام دادم
\end{itemize}

% ============================================================
\section{جمع‌بندی}

\subsection{سه نکتهٔ کلیدی}
\begin{enumerate}
    \item \textbf{Design Pattern = راه‌حل آماده}
    \begin{itemize}
        \item Creational, Structural, Behavioral
        \item الگو = قالب، نه کد
    \end{itemize}
    \item \textbf{الگو را با نیاز انتخاب کن}
    \begin{itemize}
        \item مشکل را بشناس
        \item ساده‌ترین راه‌حل
        \item از روی مد نباشد
    \end{itemize}
    \item \textbf{Design Pattern + SOLID}
    \begin{itemize}
        \item هر دو برای کد قابل نگهداری
        \item مکمل هم
    \end{itemize}
\end{enumerate}

\subsection{سؤال نهایی}
\begin{quote}
\textbf{اگر امروز کد Sprint 2 خود را ببینید، کدام قسمت را می‌شود با 
Strategy یا Observer ساده‌تر کرد؟}
\end{quote}

\vspace{1cm}

\begin{center}
\textbf{\Large پایان جزوه}\\[0.5cm]
\textbf{\Large موفق باشید!}
\end{center}

\end{document}
